Nguồn: \href{http://daidoanket.vn/giao-su-phan-dinh-dieu-nha-khoa-hoc-luon-tran-tro-voi-su-nghiep-dai-doan-ket-toan-dan-toc-435683.html}{Báo Đại Đoàn Kết}

\textbf{Hoà hợp và phát huy nguồn lực trí tuệ của dân tộc luôn là mối quan tâm trăn trở của GS Phan Đình Diệu và đã được thể hiện trong nhiều bài phát biểu tâm huyết của ông. Nhân một năm ngày ông ra đi, chúng tôi xin trích đăng những ý kiến đóng góp của ông - rút từ những trang viết "Nghĩ suy cùng đất nước”.}

Sau ngày thống nhất, tuy đất nước đã được hoà bình nhưng vẫn còn không ít bộn bề bởi cuộc chiến ngăn cách bờ cõi, và cả lòng người. “Nghĩ suy cùng đất nước”, với GS Phan Đình Diệu, trong những năm tháng còn đầy rẫy khó khăn đó của nước nhà, vì thế là mong mỏi được góp phần lấp dần những vết thương còn chưa kịp lành miệng do lịch sử để lại.

10 năm sau sự kiện lịch sử ấy, và là 1 năm trước thềm Đổi mới, trên diễn đàn Mặt trận Tổ quốc (MTTQ) năm 1985, ông vẫn không ngừng nhấn mạnh tầm quan trọng của sức mạnh đoàn kết và hòa hợp, hòa giải dân tộc mà trong đó, MTTQ đóng vai trò nòng cốt: “Nhiệm vụ của MTTQ chúng ta luôn luôn là đoàn kết dân tộc để xây dựng và bảo vệ đất nước. Vì vậy, hơn lúc nào hết, trong thời điểm hiện nay, cần tập hợp được cao nhất tất cả lực lượng của dân tộc để phấn đấu cho sự phồn vinh của đất nước. Tuy nhiên, ta cũng cần nhìn nhận rằng tuy thống nhất đất nước đã được mười năm, nhưng dân tộc và đất nước chưa thật sự được thống nhất một cách mạnh mẽ: lòng người còn phân tán, Bắc –Nam còn nhiều phân biệt, tình trạng địa phương cát cứ trầm trọng, nhiều vấn đề sau chiến tranh còn nhức nhối, còn nhiều phân biệt đối xử trong các tầng lớp nhân dân, dòng người bỏ đất nước ra đi còn tiếp diễn...

Nhưng quan trọng nhất là sức mạnh của dân tộc (về tri thức, về kinh tế, về khả năng sản xuất kinh doanh, ở trong nước và ngoài nước) chưa được huy động cho sự phát triển đất nước. Một chính sách đoàn kết dân tộc rộng rãi cần được ban bố và thi hành trong thực tế. Cốt lõi của chính sách đó phải là: mọi người Việt Nam, dù ở góc trời nào, miễn là có nguyện vọng, đều có thể tìm được một chỗ đứng bình đẳng trên đất nước Việt Nam này để lao động và cống hiến cho sự phồn vinh của đất nước bằng mọi cách thích hợp với hoàn cảnh và khả năng của mình. Và tất nhiên, điều đó sẽ được thể hiện bằng những chính sách cụ thể, như về giải quyết các hậu quả sau chiến tranh còn tồn tại, chính sách đối với người Việt Nam ở nước ngoài, chính sách thu hút vốn đầu tư và thu hút lực lượng khoa học kĩ thuật...”. Những ý kiến này của ông là một phần trong bài phát biểu “Góp ý kiến về đổi mới” mà toàn văn đã được đăng 3 năm sau đó trên báo Tổ Quốc năm 1988.

20 năm sau, cũng trên diễn đàn MTTQ, trong bài phát biểu “Vấn đề đoàn kết dân tộc dưới cách nhìn mới của khoa học”, sau khi trình bày về “tư duy hệ thống” trong khoa học về sự phức tạp (khoa học mà nhà vật lý nổi tiếng Stephen Hawking đã nhận định sẽ là khoa học của thế kỷ 21), ông nói: “Đất nước ta, thế giới mà ta đang sống, đều là những hệ thống thích nghi phức tạp. Từ sau khi thống nhất, đất nước ta đã có nhiều thay đổi. Về cơ bản, trên đất nước ta, những mâu thuẫn đối kháng một mất một còn phải được giải quyết bằng đấu tranh ai thắng ai đã không còn, một mục tiêu chung vì dân giàu nước mạnh, xã hội công bằng dân chủ văn minh đã trở thành cái lẽ chung cho sự đồng thuận xã hội. Tất nhiên nói như thế không có nghĩa là đất nước đã ở trong một trạng thái cân bằng ổn định, vả chăng cái cân bằng ổn định cũng đồng thời là cái trì trệ không đổi là cái mà ta không hề mong đợi. Vẫn còn đó nhiều cái khác nhau, thậm chí đối chọi và cạnh tranh nhau, cùng chung sống với nhau và tương tác với nhau trong những quan hệ hết sức đa dạng. Đoàn kết toàn dân tộc ngày nay không còn có ý nghĩa là đòi hỏi phải thống nhất một ý chí và hành động theo một sự chỉ đạo duy nhất, mà phải là trong môi trường của sự đồng thuận chung, mọi thành phần được tự do phát huy mọi năng lực trong hoạt động của mình, để rồi qua sự tương tác hết sức đa dạng của xã hội tạo nên những hợp trội, tức là những chất lượng mới, những trật tự mới làm nên những giàu có chung của toàn xã hội...”. 

"Làm giàu cho đất nước” được bắt nguồn từ sự “làm giàu” những cảm xúc công dân, khiến họ cảm thấy “được hưởng niềm ấm áp xem đất nước này là của mình”, để “mọi thành viên được cống hiến cho đất nước phần tươi mát và trong lành nhất của trí tuệ và sức lực”... (trích một bài viết của ông trên báo Sài Gòn Giải Phóng, 1989).

Trong bức thư gửi nguyên Tổng Bí thư Đỗ Mười, ngày 10/8/1991, thêm lần nữa, ông lại bày tỏ mối quan tâm đau đáu ấy của mình: “Trong nội bộ dân tộc ta, ước gì ta có thể xem tất cả mọi người đều là bà con, anh em, bạn bè, có thể cùng sống chung trên mảnh đất thân yêu này trong đoàn kết và hòa hợp. Lợi ích của dân tộc, sự phồn vinh của đất nước đòi hỏi vượt qua chia rẽ và thù hận, cả cái không tránh được và cái không đáng có, để tập hợp đến tối đa mọi nguồn lực lành mạnh của dân tộc cùng phấn đấu cho sự giàu mạnh của Tổ Quốc (...). Đoàn kết trên cơ sở thật sự phát huy dân chủ và tập hợp trí tuệ của dân tộc sẽ là cách tốt nhất huy động mọi lực lượng cùng xây dựng đất nước trong sự ổn định của những thay đổi cần thiết được chủ động trù tính trước. Với mục đích đó, nên khuyến khích các hình thức tiếp xúc, trao đổi, đối thoại giữa các tầng lớp nhân dân, trong nước và ngoài nước, giải tỏa những mặc cảm căng thẳng không đáng có, tạo không khí hòa hợp rất cần cho sự phát triển trong ổn định hiện nay...”

Tại hội nghị của UBTƯ MTTQ Việt Nam tháng 3/1993, vị Ủy viên Đoàn chủ tịch đã nhấn mạnh về tính tối cần của việc huy động năng lực trí tuệ trong và ngoài nước nhằm giúp đất nước tiến lên: “Ngoài những năng lực hiện có ở trong nước chưa được hoặc chưa có điều kiện để phát huy, chúng ta còn những năng lực trí tuệ của một số bộ phận đông đảo người Việt Nam ở nước ngoài. Những người này hiện đang hướng về đất nước với nhiều nhiệt tình, nhưng cũng với nhiều nghi ngại. Làm thế nào để họ bỏ nghi ngại, phát huy được năng lực của mình thì đó sẽ là một năng lực rất to lớn của đất nước. Trí tuệ là của cải. Trí tuệ là một tài sản còn quan trọng hơn tiền bạc rất nhiều...”.

Đây cũng chính là mối quan tâm của ông từ hơn 10 năm trước đó, trong chuyến đi kéo dài 6 tuần (tháng 9-10/1980) tới 15 thành phố lớn và 14 trường Đại học Mỹ, 3 trường ở Canada và một số cơ quan nghiên cứu ở Pháp..., cũng là dịp để ông có được những cuộc gặp thân tình với những người đồng hương của mình tại xứ người: “Một đêm đọc thơ, ngâm thơ, bình thơ thật đáng nhớ. Ôi! Một “quê hương có con sông xanh biếc” làm da diết biết bao nhiêu tấm lòng của những kẻ tha phương! Tình quê, vâng, tình quê, phải chăng đó là sợi dây thần mãi mãi gắn bó mọi tâm hồn dân Việt, dù họ đến tự nơi nào, và họ sẽ đi đâu, về đâu...”. 

Và cùng đó, là những suy tư mong đợi: “Tôi vẫn không rời khỏi được những ý nghĩ về khoảng 300.000 người Việt Nam hiện có trên đất Mỹ. Phải chăng phần đông trong số họ cũng nhớ đất nước, quê nhà? Và rồi mười năm sau, hai mươi năm sau, với tinh thần hiếu học, cần cù của người Việt Nam nói chung, lại được nền khoa học kỹ thuật tiên tiến ở đây dìu dắt, hẳn là sẽ có hàng chục ngàn nhà kỹ thuật, kinh tế và khoa học người gốc Việt có tài năng trên đất này. Có cách gì không nhỉ, để rồi sau mươi năm, hai mươi năm, đa số những chuyên gia người Việt đó sẽ nhìn về Tổ Quốc với những mong muốn đóng góp của những đứa con xa nước xa nhà?” (Trích “Nhật ký một chuyến đi xa” của GS Phan Đình Diệu). Như chính những câu thơ được ông viết trước đó, “từ Cuba xa xôi”: “Khi đất Á, khi trời Âu biển Mỹ/ Vẫn dõi lòng người qua những đại dương/ Ta hiểu tình ta từ nửa vòng trái đất/ Hiểu cái lắng sâu của cuộc sống bình thường / Bởi tự rất xa nhìn cái gần mới thật/ Mới rõ chiều dài của năm tháng yêu thương/ Ôi xa thẳm là mênh mang trời bể/ Hay tình người từ muôn dặm trùng dương...”.

Tại nơi an nghỉ của nhà khoa học nổi tiếng có tâm, có tầm, luôn trăn trở “nghĩ suy cùng đất nước”, thêm lần nữa, chúng ta lại được gặp lại câu nói khiêm nhường, rất Phan Đình Diệu: “Tôi mong là một giọt nước nhỏ bé hòa vào nhiều triệu giọt nước khác của dân tộc để tạo thành dòng thác đổi mới cho đất nước”...


\textit{GS Phan Đình Diệu (1936-2018) - Nhà toán học, nhà khoa học máy tính nổi tiếng, được ghi nhận là người đặt viên gạch đầu cho sự phát triển của ngành tin học VN. GS là người sáng lập và là Chủ tịch đầu tiên của Hội Tin học Việt Nam; đồng thời là Viện trưởng đầu tiên của Viện Khoa học tính toán và Điều khiển (nay là Viện Công nghệ Thông tin Việt Nam)... Trong công tác xã hội, ông là Ủy viên Đoàn chủ tịch Uỷ ban Trung ương Mặt trận Tổ quốc Việt Nam các khoá III, IV, V, VI, VII; nguyên Đại biểu Quốc hội Việt Nam khóa V, VI... Ông được biết là người luôn có những phát biểu thẳng thắn và tâm huyết nhằm đóng góp vào sự đổi mới của đất nước. }